% Response to reviewer and editor comments
% JAWRA MS 8299185 -- "A real-time reservoir inflow forecast evaluation framework"
%
% Standalone document; compiles with pdflatex. No external .bib needed.

\documentclass[11pt]{article}

\usepackage[margin=1in]{geometry}
\usepackage{amsmath}
\usepackage{booktabs}
\usepackage{xcolor}
\usepackage{enumitem}
\usepackage{titlesec}
\usepackage{hyperref}
\usepackage{parskip}

\hypersetup{colorlinks=true, linkcolor=black, urlcolor=blue!60!black}

% Allow long URLs to break across lines.
\usepackage{xurl}

\definecolor{quotegray}{gray}{0.32}

% Reviewer comments are set as indented gray italics.
\newenvironment{rc}%
  {\begin{list}{}{\setlength{\leftmargin}{1.5em}\setlength{\rightmargin}{1.5em}}%
   \item[]\itshape\color{quotegray}}%
  {\end{list}}

% New or revised manuscript text.
\newenvironment{newtext}%
  {\begin{list}{}{\setlength{\leftmargin}{1.5em}\setlength{\rightmargin}{1.5em}}%
   \item[]}%
  {\end{list}}

\titleformat{\section}{\large\bfseries}{\thesection}{0.6em}{}
\titleformat{\subsection}{\normalsize\bfseries}{\thesubsection}{0.6em}{}

\setlength{\parskip}{0.6em}

\title{\vspace{-2em}Response to reviewer and editor comments}
\author{}
\date{}

\begin{document}

\maketitle
\vspace{-4em}

\noindent\textbf{Manuscript:} ``A real-time reservoir inflow forecast evaluation
framework'' (MS 8299185), submitted to the \emph{Journal of the American Water
Resources Association}.

\noindent\textbf{Authors:} Cameron Bracken, Youngjun Son, Vince Tidwell, Nathalie Voisin.

\vspace{1em}

We thank the editor and the four reviewers for their comments. The reviewers and the
handling editor converged on one point: the manuscript described a framework without
stating what was new about it. We agree. The original draft presented the five-step
process as though the process itself were the contribution. We have added an explicit
statement of contribution to the introduction, restructured the process description so
that the generic steps are labeled as such, added a discussion subsection comparing the
framework against retrospective verification, and added a result in which a tailored
metric separates two forecast products that the generic metrics do not.

We have also expanded the description of the study system and the forecast products,
added a smoothing sensitivity analysis, characterized the hydrology of the evaluation
period, cited the forecast value literature, and corrected several errors in the methods
section that the reviewers identified.

Two statements in the original manuscript were wrong and we correct them here:

\begin{itemize}[itemsep=0.15em]
    \item The evaluation period was given as July 2023 to May 2025. The record analyzed
    runs from June 2023 to April 2025.
    \item The three forecast products do not cover the period uniformly. Forecast B spans
    the full period; forecast A and the in-house forecast both end in November 2024. The
    metrics in Table 1 are therefore computed over overlapping but not identical samples.
    We now report the sample size for each product and have verified that restricting the
    comparison to the 510 days on which all three are present does not change the ranking
    at any location.
\end{itemize}

Reviewer comments are set in indented italics. New or revised manuscript text is set
indented in roman type.

\vspace{1em}
\hrule
\vspace{1em}

\section*{Editor comments}

\begin{rc}
The greatest barrier, from my reading of the manuscript and review of the comments,
is the question of novelty and reproducibility. Reviewers 1, 3, and 4 note that
there is no distinct statement of novelty, with reviewer 3 particularly noting that
the manuscript seems to suggest that the framework is the novelty without fully
describing the novelty.
\end{rc}

Agreed. We have added a statement of contribution to the end of the introduction:

\begin{newtext}
``The contribution of this work is a specific operating procedure for the case where a
utility must choose between competing forecast products under live operating conditions
and cannot run a hindcast. That procedure has two components that distinguish it from
retrospective verification practice. First, the operators participate in metric selection
and design and can modify, remove, or append metrics to the evaluation as more data is
collected. Second, the evaluation runs on the same operational data the operators use
which reflects the real day-to-day operational decisions.''
\end{newtext}

We have also rewritten the preceding paragraph so that the gap is stated directly rather
than implied. Traditional verification requires hindcasts over a climatologically
representative period, which is expensive and often infeasible. The fallback of evaluating
archived operational forecasts fits poorly with products that are updated continuously,
trained on evolving data streams, hand adjusted by operators, or developed by third
parties.

The claim that operator participation changes the outcome is now supported by a result
rather than asserted. The forebay exceedance metric exists only because the operators
asked for it, and it separates the two commercial products where the generic metrics do
not. From a neutral forebay start all exceedance probabilities are below 10\% and the
products are difficult to distinguish. From the non-neutral start operators are routinely
in, forecast A reaches 23.7\% at Bellows Falls against 1.8\% for forecast B. This is now
presented as a result in the tailored metrics section rather than left implicit.

On reproducibility, the evaluation code is public
(\url{https://github.com/HydroWIRES-PNNL/inflow-forecast-evaluation}) and sample data are
on Zenodo (\url{https://doi.org/10.5281/zenodo.16921728}). We have added the parameter
values that were previously left implicit, most importantly the smoothing window
(comment 4.1), together with a sensitivity analysis over it; the drainage areas and USGS
reference gages for all six evaluation points (comment 1.5); the issue frequency and
horizon of each product; and the per-product sample sizes. The proprietary forecast time
series cannot be released, which we now state as a limitation.

\vspace{1em}
\hrule
\vspace{1em}

\section*{Reviewer 1}

\begin{rc}
The topic and the contents are relevant, and it is very interest for hydropower
operators and decision-makers. However, the novelty and main contributions of this
study are unclear.
\end{rc}

Thank you for the review. The comments on framing, the missing system description, and the
table and figure quality were all warranted. Your comment 1.6, asking us to check an
apparent anomaly in the NSE values, identified something informative rather than an error,
and it has become a result in the revised manuscript.

\subsection*{1.1 Novelty and main contributions}

\begin{rc}
The novelty and main contributions of this study are unclear. Please explain and
discuss in detail.
\end{rc}

Please see our response to the editor for the new statement of contribution. In brief, the
contribution is a procedure for the case where a utility must choose among competing
forecast products under live operating conditions and a hindcast is not available. What
distinguishes it from standard verification is that metric selection is driven by operator
input rather than fixed in advance, and that the evaluation runs on the operational data
feed rather than an archived dataset.

We have also added to the discussion why this matters now. Forecast products using AI and
machine learning methods are being marketed to utilities at an increasing rate, and the
verification statistics vendors supply may not be verifiable. A utility generally has no
standard procedure for checking those claims against its own system, and as we now report,
often cannot inspect the model at all.

\subsection*{1.2 Persistence as a benchmark at long lead times}

\begin{rc}
Lines 70--72: ```persistence' forecasts which use the last observed value carried out
for all future timesteps to get a sense of the arguably worst possible forecast.
Persistence is a competitive benchmark for short term operations, from hours to 2-day
horizons''. However, Fig.\ 3 and Fig.\ 4 show that 10-day and months ``persistence
forecasts'' as benchmark which is unrealistic in practice.
\end{rc}

You are correct, and this was an internal inconsistency: we stated that persistence is
competitive only to about two days and then plotted it out to ten. Reviewer 2 raised a
related point in comment 2.7. We have revised the benchmark section:

\begin{newtext}
``Persistence is a simple but competitive forecast at short lead times, from hours to
roughly two days, and is the appropriate reference for day-ahead operations, which is the
horizon of interest for the system evaluated here. It degrades quickly beyond that horizon,
where a climatological forecast is the more appropriate lower bound. Evaluators should
match the benchmark to the horizon of interest: persistence for hourly to 2-day operations,
climatology for weekly to seasonal horizons. Where persistence is plotted at longer lead
times in the figures below, it is retained only to show the scale of degradation and should
not be read as a meaningful reference at those horizons.''
\end{newtext}

The figure captions now carry the same qualification.

We have not added a climatological benchmark to the figures. A defensible daily climatology
for these locations requires a long record of calculated inflow at the dams, and that
record begins with the evaluation. A climatology derived from the nearby USGS gages would
differ in drainage area and regulation, so it would not be commensurable with the quantity
being verified. We have stated the recommendation, which is general, and been explicit that
the day-ahead horizon is where persistence is the correct reference for this system.

\subsection*{1.3 Operating problems and the forecasting scheme in practice}

\begin{rc}
What are the main operation problems of these hydropower stations? Which forecasting
scheme is used in practice?
\end{rc}

This context was missing and is needed to make the tailored metrics comprehensible. We have
added a subsection, ``Operating context.'' The three dams have limited storage relative to
their inflow and are operated as run-of-river facilities, with a daily cycle of an upward
ramp, flexible generation, a downward ramp, and refill. Generation is bid into the ISO New
England day-ahead market, so the day-ahead inflow volume is the quantity that matters most.
The binding constraint comes from the FERC relicensing of the three lower Connecticut River
dams, which imposed a $\pm$0.5 foot ($\approx$15~cm) limit on forebay fluctuation with
penalties for exceedance. The central operating problem is committing generation a day
ahead on an uncertain inflow forecast without violating the forebay limits, which is what
motivated the cumulative-error and threshold-exceedance metrics.

On the scheme used in practice: the operators used an in-house forecast built from National
Weather Service River Forecast Center forecasts, routed to the dam locations and adjusted
using operator judgment. This is the incumbent product against which the two commercial
forecasts were evaluated, labeled GRH in the results. It extends only about 32 hours,
because the day-ahead bid is the only decision it is used for.

\subsection*{1.4 Description of the forecast products}

\begin{rc}
Two commercial forecasts (A, B) and the GRH in-house forecast (GRH) should be
described in detail, such as what type of hydrologic model used? How about the
calibration and validation results, what type input data requirement etc.
\end{rc}

Reviewers 2 and 4 raised the same point (comments 2.16 and 4.3). We agree that
anonymization without any methodological characterization leaves readers unable to
interpret the results. We have added a subsection describing what we are able to disclose:

\begin{newtext}
``Forecast A is produced by a consulting firm using a conceptual rainfall-runoff model with
hydraulic routing through the Connecticut River mainstem, driven by a commercial numerical
weather prediction forecast, and issued once daily at an hourly timestep out to roughly 8
days. Forecast B is produced by a technology firm using a machine learning model trained on
gauge and remotely sensed catchment observations, issued twice daily at hourly timestep out
to 10 days. The GRH in-house forecast combines RFC forecasts with routing and operator
judgment and extends only through the day-ahead period, roughly 32 hours, because that is
the horizon it is used for. Neither vendor released model internals or calibration and
validation diagnostics, and we did not have access to them. A utility evaluating commercial
forecast products generally cannot inspect them and must judge them on operational
performance, which is part of the motivation for this framework.''
\end{newtext}

The issue frequencies and horizons are taken from the forecast archive rather than vendor
documentation, so they describe what was actually delivered.

We cannot provide calibration and validation results. Our agreements prevent us from naming
the vendors or describing model internals, and the vendors did not release diagnostics to
us. We have chosen to state this plainly and treat it as a finding about the procurement
situation operators face. We have also added that a full procurement evaluation should weigh
operational cost, data dependencies, and ease of integration alongside skill, since these
differ substantially between the two commercial products.

\subsection*{1.5 Watershed information}

\begin{rc}
Three main hydropower producing dams in the system, Wilder, Bellows Falls, and Vernon
(Figure 2) as well as three main tributaries, the Ottauquechee, Sugar and White
Rivers. The watershed information, such as area, river length etc.\ should be
described in detail.
\end{rc}

Agreed. We have added a table of physical characteristics for all six evaluation points,
with drainage areas from the USGS National Water Information System for the reference gage
at each location:

\begin{center}
\begin{tabular}{lllrr}
\toprule
Location & Type & USGS reference gage & Area (mi$^2$) & Area (km$^2$) \\
\midrule
Wilder             & Mainstem dam & 01144500 West Lebanon, NH   & 4,092 & 10,598 \\
Bellows Falls      & Mainstem dam & 01154500 North Walpole, NH  & 5,493 & 14,227 \\
Vernon             & Mainstem dam & 01156500 Vernon, VT         & 6,266 & 16,229 \\
White River        & Tributary    & 01144000 West Hartford, VT  &   690 &  1,787 \\
Sugar River        & Tributary    & 01152500 West Claremont, NH &   269 &    697 \\
Ottauquechee River & Tributary    & 01151500 North Hartland, VT &   221 &    572 \\
\bottomrule
\end{tabular}
\end{center}

We have also added text describing the basin. The mainstem drainage areas are nested, so
the area above Vernon includes those above Bellows Falls and Wilder. This nesting explains
the downstream propagation of bias noted in the original manuscript, which we now quantify:
for forecast A, day-ahead PBIAS at Bellows Falls is 0.4\% and grows to $-19.1$\% by day 7,
with larger growth downstream. Because releases from each dam become inflow to the next, an
error in an upstream forecast propagates to the downstream forecast and compounds with lead
time.

Basin hydrology is snowmelt influenced, with a spring freshet typically spanning March
through May, convective rainfall in summer, and rain and rain-on-snow driven events in
winter. Several tributaries have upstream flood control regulation by the U.S.\ Army Corps
of Engineers. We have not added river lengths, since the drainage areas and the nesting
relationship are what the results depend on.

\subsection*{1.6 NSE values in Table 1}

\begin{rc}
In Table 1, the NSE values of forecast scheme A at Wilder, Bellows and Vermon
watersheds are 0.72, 0.55, 0.81, respectively. Please check the evaluation results.
The forecasting accuracy at the Bellows might high than that at the Wilder.
\end{rc}

Thank you for checking this. We re-ran the calculation and the values are correct as
published. Your question prompted us to look into why, and the answer is a normalization
effect that we had not explained, so we have added an explanation rather than only
confirming the numbers.

NSE and KGE are normalized by the variance of the observations at each location, so their
values are not directly comparable across locations with different inflow variability
(Knoben et al., 2019). Flows become progressively more regulated and
therefore smoother downstream, and regulated flows have less variability, so the more
regulated a location the higher its NSE tends to be. The persistence benchmark makes this
visible: persistence alone achieves NSE of $-0.81$ at Wilder, 0.39 at Bellows Falls, and
0.80 at Vernon. At Vernon persistence already explains most of the variance, leaving little
room for a forecast to improve on it.

We have added a KGE skill score column to Table 1, computed against the persistence
forecast, which removes the normalization problem. It reverses the apparent ordering.
Forecast A is weakest at Vernon (skill score 0.47), not at Bellows Falls (0.65), and its
skill at Vernon is well below the other two products there (0.81 for forecast B and 0.60
for the in-house forecast), a gap that the raw KGE values of 0.83, 0.94, and 0.88
substantially understate.

So the direct answer to your question is that the cross-site pattern in the raw NSE values
is largely a normalization artifact and should not be read as a ranking of forecast quality
across locations. We have added this to the manuscript as an illustration of why single-metric
interpretation is hazardous:

\begin{newtext}
``This example shows how interpretations based on a single metric can lead to potentially
erroneous conclusions and underscores the importance of using multiple metrics in any
evaluation, especially when comparing forecasts between locations.''
\end{newtext}

We also now recommend in the methods section that skill scores relative to a benchmark be
reported alongside raw metrics in any multi-site evaluation. This was a real gap and we are
grateful the question was raised.

\subsection*{1.7 Table format}

\begin{rc}
The format of Tables in the manuscript are poorly designed and needed be changed.
\end{rc}

Agreed. Table 1 has been rebuilt with rows grouped by location, the location named once per
group, numeric columns aligned, consistent significant figures across metrics, and units in
the column headers. It has gained three columns: the number of day-ahead forecasts evaluated
for each product, which is necessary because the products cover different portions of the
period; RMSE in cfs, requested in comment 4.7; and a KGE skill score against the persistence
forecast, described in comment 1.6. Two new tables have been added, one giving basin
characteristics (comment 1.5) and one giving the smoothing sensitivity analysis (comment
4.1), and a third traces each tailored metric to the operator concern that prompted it
(comment 3.2).

\subsection*{1.8 Figures}

\begin{rc}
The figures also need be redrawn.
\end{rc}

We have revised the captions substantially and describe below what we have and have not
changed, since we would rather be precise than claim a wholesale redraw.

Captions are now self-contained. Every caption states the quantity plotted, the meaning of
any reference line, and the units. The Figures 3 and 4 captions now give the no-skill
threshold values, NSE of 0 and KGE of $-0.41$, rather than referring vaguely to ``no skill
over the mean,'' and note that persistence is a meaningful reference only over the first two
days (comment 1.2). The Figure 5 caption specifies that the error bars span the 2.5th to
97.5th percentile (comment 4.8) and states the pass-inflow assumption and the compliance
band. The Figure 1 caption states the smoothing window and gives the SI conversion for cfs
(comment 2.17). The Figure 6 caption explains why the lower panel is the operationally
relevant case.

On the figures themselves we would welcome more specific guidance. We are willing to change
the palette, line styles, text sizes, and output format, and we agree the legibility of axis
and legend text at print size could be improved. We have not regenerated them because we did
not want to redraw them on a guess at the objection and then again after learning what it
was. If the reviewer would indicate which aspects they consider deficient we will address
them directly.

\vspace{1em}
\hrule
\vspace{1em}

\section*{Reviewer 2}

\begin{rc}
The core work is solid and the practical contribution is genuine. Assuming the major
comments are addressed I recommend this manuscript is accepted with minor revision.
\end{rc}

Thank you for a detailed and technically careful review. Several of the minor comments
identify outright errors in the methods section, including one in an equation, and we are
grateful they were caught. The two major comments have both changed the manuscript
substantially.

\subsection*{2.1 Engagement with the forecast value literature}

\begin{rc}
The stated goal of the framework is to help operators evaluate whether a forecast
product will improve operations enough to justify investment. The paper does not
include any assessment of forecast value and the authors acknowledge this gap
(L172--177), framing it as ``an important next step.'' [\ldots] The authors do not need
to perform a full value assessment, but they should reference this literature, discuss
how it connects to their framework, and provide operators with at least a conceptual
pathway from the metrics presented to the investment decision the paper is motivated by.
\end{rc}

This is a fair criticism and identifies a genuine gap between what the introduction promises
and what the paper delivers. We have added a discussion subsection, ``From forecast quality
to forecast value'':

\begin{newtext}
``The metrics presented here are measures of forecast quality, not forecast value. Quality
is a statistical property of a forecast paired with observational data. Value depends
additionally on human decisions and the cost associated with forecast errors. A forecast can
be more statistically accurate and yet carry less value, if the additional accuracy falls in
a range where the operator would not act differently. We refer the reader to Farkas et al.\
(2025) for a comprehensive review of the economic value of forecasts.

The tailored metrics are a step towards assessing forecast value. For example, the forebay
exceedance probability is a quantity with a direct financial consequence, since exceedance
incurs a penalty under the FERC license. This is a partial valuation, since it captures an
upper limit on compliance risk but it is not put in terms of generation revenue. A more
complete assessment of revenue requires a decision-support model which includes reservoir
operations and market conditions. We recommend that operators treat a real-time quality
evaluation as the first of two stages. The evaluation establishes which forecast products
are plausibly worth adopting and produces the archive of forecasts that may be used to drive
the second step where value is assessed.''
\end{newtext}

We have cited Farkas et al.\ (2025), a recent systematic review, in preference to citing
individual methods. We drafted a more detailed treatment of specific value metrics, including
the relative utility value approach, but judged that a paragraph comparing value metrics we
do not apply would read as padding in a paper about quality evaluation. If you or the editor
would prefer the fuller treatment we are glad to restore it.

We have not attempted a value assessment. Doing so credibly requires the utility's damage
function and a representation of how operators respond to the forecast, which returns to the
scheduling model coupling discussed under comment 4.9.

\subsection*{2.2 Positioning of the contribution}

\begin{rc}
The five-step iterative process (collect data, compute metrics, visualise, solicit
feedback, iterate) reads more as standard good operational practice than a novel
framework contribution. I feel that the tailored metrics and operator feedback loop
are where the genuine contribution lies, especially in the real-time context, and
these deserve more prominence and deeper treatment.
\end{rc}

We agree, and your diagnosis matches Reviewer 3's and the handling editor's. Yours says not
only what is missing but where the contribution actually is, and we have restructured
accordingly. Please see our response to the editor for the new statement of contribution.

The specific changes: steps 1 through 3 are now labeled as generic practice for which we
claim no novelty. Soliciting feedback and acting on it have been separated into two steps,
``Gather feedback'' and ``Update metrics,'' since collecting operator comment and changing
the metric set are different activities and the second is the one that matters; the process
is now six steps rather than five. The mechanics of both are documented rather than asserted,
including a table tracing each tailored metric to the operator concern that prompted it
(comment 3.2). The two timescales conflated in the iteration step, metric convergence and
record length, are now separated (comment 3.3). We have added the case where a tailored
metric separates two products that the generic metrics do not (comment 3.5), and a discussion
subsection comparing the framework against retrospective evaluation (comment 3.4). We have
also added that operational reliability is part of what the evaluation measures: a forecast
that cannot be produced operationally should not be considered in a real-time evaluation.

\subsection*{2.3 Hyphenation of ``real time''}

\begin{rc}
(L2, Title) ``real time'' should (perhaps) be hyphenated as ``real-time'' when used as
an adjective (also applies throughout the manuscript, e.g.\ L44, L48, L50, L110, L116).
\end{rc}

Corrected throughout, including the title, which is now ``A real-time reservoir inflow
forecast evaluation framework.'' Reviewer 4 raised the same point in comment 4.5. The
manuscript contained 29 instances of ``real time'' and none of the hyphenated form. We
reviewed each and hyphenated in attributive position (``real-time evaluation''), leaving the
unhyphenated form where it functions as an adverbial phrase (``used in real time'').

\subsection*{2.4 Verification improving alongside forecasts}

\begin{rc}
(L12--13) The sentences seem to imply that we should expect verification methods to
improve alongside improvements in forecasts, but this isn't the case. Also, the first
sentence is about forecast and scheduling methods but the second sentence links it to
operational paradigms somewhat abruptly. Suggest restructuring so the research gap is
clearer.
\end{rc}

You are right on both counts and the implication was not intended. We have rewritten the
passage so the gap is the subject:

\begin{newtext}
``The quality of hydrologic forecasts and hydropower scheduling has rapidly improved in
recent years due to the emergence of new data driven approaches and scheduling optimization.
Despite this, the practice of forecast verification has remained mostly unchanged. A
traditional forecast verification requires evaluation of hindcasts using standard statistical
metrics, over a period which is representative of the full climatological record, often 30
years or more. The high cost of hindcasts and strict data requirements are two reasons full
verifications are not done in practice. The second best approach, if hindcasts are not
available, is to evaluate archived operational forecasts, although this approach fits poorly
with products that are updated continuously, trained on evolving data streams, hand adjusted
by operators, or developed by third parties. The gap this study addresses is to develop a
procedure by which hydropower or reservoir operators can evaluate forecasts on their own
system under familiar operating conditions.''
\end{newtext}

The point about products delivered as services rather than inspectable models is now
supported later in the paper, where we report that neither vendor released model internals or
calibration diagnostics.

\subsection*{2.5 Dismissal of ensemble metrics}

\begin{rc}
(L42--44) The distinction between deterministic and probabilistic is important and
well noted. However, the dismissal of ensemble evaluation metrics seems premature
given that the authors later suggest probabilistic forecasts may help (L170). Consider
briefly noting what ensemble metrics could be added if ensemble forecasts were
available.
\end{rc}

Agreed, and the inconsistency is a fair catch: we set probabilistic evaluation aside early
and then invoked probabilistic forecasts as a remedy later. Reviewer 4 asked for the same
addition in comment 4.4, so we have added a methods subsection, ``Extension to probabilistic
and ensemble forecasts,'' covering CRPS, rank histograms, and reliability diagrams. Please
see comment 4.4.

We have also softened the framing so that the focus on deterministic forecasts is presented
as a property of the products under evaluation rather than a judgment about ensemble metrics.
The revised text notes that all three products evaluated here are deterministic, which is why
the metrics we describe are deterministic, and points forward to the new subsection.

\subsection*{2.6 The argument against hindcast evaluation}

\begin{rc}
(L45--47) The argument against hindcast evaluation is reasonable but somewhat
overstated. Hindcast evaluations do not necessarily require operators to re-operate
the system, especially if verifying forecasts of environmental variables such as
inflows, which you are. Suggest softening this claim or explaining how operators are in
the loop for the inflow forecasts.
\end{rc}

You are correct and the claim as written was overstated. Verifying an inflow forecast against
observed inflow does not require re-operating the system. The re-operation requirement applies
to the value assessment rather than the quality assessment, which is the distinction your
first major comment prompted us to draw. We have rewritten the passage:

\begin{newtext}
``A hindcast evaluation of forecast quality does not itself require re-operating the system,
since forecasts of inflow can be compared directly against observed or calculated inflow.
Yet, two complications arise in this setting. First, where the forecast is modified using
human judgment, a hindcast requires forecasters to reconstruct the decisions they would have
made under past conditions. Those conditions may include electricity prices, weather
conditions, or other operational conditions that may be difficult to reconstruct after the
fact. Second, extending the evaluation from forecast quality to forecast value (eg.
streamflow to megawatts) requires a reservoir operations model or manual operator input. A
real-time evaluation avoids both complications by developing a forecast archive over time and
relying on real operator decisions.''
\end{newtext}

There is a further dependency worth stating, which we had not made explicit. The calculated
inflows we verify against are derived from forebay elevation changes and dam outflow, so the
observation record is itself a function of how the system was operated. This does not require
re-operation to interpret, but it does mean our reference data is not independent of
operations in the way a gauge record would be. We have added this to the limitations.

\subsection*{2.7 Persistence as ``arguably worst possible forecast''}

\begin{rc}
(L71) Describing persistence as ``arguably worst possible forecast'' is incorrect.
Persistence is a competitive benchmark for short lead times as the authors themselves
note in the next sentence. Climatology is typically considered the lower bound for
longer horizons. Suggest rephrasing to ``a simple but competitive benchmark.''
\end{rc}

Corrected using your suggested wording. Reviewer 1 raised the related point about plotting
persistence at long lead times in comment 1.2; that response includes the revised passage and
explains why we have not added a climatological benchmark despite recommending one.

Our own results make your point concretely: persistence achieves day-ahead NSE of 0.80 at
Vernon, which is competitive with the commercial products, and this is what reveals that
forecast A adds little skill there.

\subsection*{2.8 Smoothing window width}

\begin{rc}
(L80--81) A triangular smoother is reasonable but the choice of window width is not
discussed. The smoothing window will affect the comparison with forecasts and should be
selected with care. Please note the window used and perhaps the sensitivity to this
choice on your findings.
\end{rc}

Reviewer 4 raised this as their first major comment and we have addressed it with both the
value and a sensitivity analysis. The window is 24 hours, implemented as two successive
24-hour moving averages of the hourly calculated inflow, which produces the triangular
weighting. It was chosen to match the day-ahead operational timescale. Metrics improve
modestly with window width, the largest change being 0.05 NSE, but the ranking of the three
products is unchanged at every window width at every location. Please see comment 4.1 for the
table.

\subsection*{2.9 KGE name and citation}

\begin{rc}
(L92) ``King-Gupta Efficiency'' should be ``Kling-Gupta Efficiency.'' Also, ``(refs)''
is a placeholder that needs to be replaced with actual references (Gupta et al.\ 2009
presumably).
\end{rc}

Both corrected, with Gupta et al.\ (2009) as you surmised, and Nash and Sutcliffe (1970)
added for NSE in the same sentence. Reviewer 4 flagged the same two items in comment 4.6.
Thank you for catching the misspelling.

\subsection*{2.10 The claim about KGE and large values}

\begin{rc}
(L92--93) The claim that KGE is not skewed by large values while NSE and RMSE are
requires clarification. KGE and NSE both use correlation as a component. The relevant
difference is that KGE decomposes performance into correlation, variability ratio, and
bias ratio rather than using squared error. The statement as written is misleading.
\end{rc}

You are right, and your characterization is the correct one. The original sentence attributed
the difference to the squared error term, but KGE also incorporates correlation, which is
itself computed from squared deviations, so the stated reason does not hold. We have replaced
it with your framing:

\begin{newtext}
``We recommend KGE because it decomposes forecast performance into correlation, a variability
ratio, and a bias ratio, so that a low score can be attributed to a specific deficiency
rather than only registered as large aggregate error. NSE and RMSE aggregate squared error
into a single number, which makes them more strongly influenced by the largest errors and
less diagnostic about their source.''
\end{newtext}

\subsection*{2.11 Placement of the Knoben et al.\ citation}

\begin{rc}
(L95--96) The KGE threshold of $-0.41$ is attributed to comparing against the mean of
observations but the Knoben et al.\ (2019) citation is on L97, not adjacent to the claim.
Consider restructuring so the citation directly supports the threshold statement.
\end{rc}

Corrected. The citation now appears with the threshold statement it supports, and we have made
the connection to the following sentence explicit:

\begin{newtext}
``Values of KGE that are greater than $-0.41$ are better than using the mean of the
observations as a predictor (Knoben et al., 2019), but larger values are preferred up to the
maximum value of 1 for a perfect forecast. Because this threshold is not zero, NSE and KGE
have different properties and should not be used interchangeably, and intuition built on NSE
values does not transfer to KGE.''
\end{newtext}

Knoben et al.\ (2019) is also directly relevant to the cross-location comparability point
Reviewer 1 raised in comment 1.6, so we now cite it in the benchmark section as well.

\subsection*{2.12 Bias versus error}

\begin{rc}
(L98--99) ``The bias of a single forecast point is the difference between observed and
forecasted values'' defines error (sometimes called residual), not bias. Bias is a
systematic tendency. Suggest correcting the definition.
\end{rc}

Corrected. The text now reads:

\begin{newtext}
``The error, or residual, of a single forecast is the difference between the forecast and the
observation. Bias is the systematic component of error, the tendency of a forecast to be too
high or too low on average, and is conventionally expressed for a set of forecasts as a
percentage of total observed volume.''
\end{newtext}

The original definition also had the sign backwards relative to the PBIAS equation that
follows it, written as $\sum(f_i - o_i) / \sum o_i$. The revised wording is consistent with
the equation.

\subsection*{2.13 Pearson correlation formula}

\begin{rc}
(Eq.\ 3) The Pearson correlation coefficient formula is missing a summation over the
second term of denominator.
\end{rc}

Corrected, and thank you for checking the equations. The denominator was written as a single
summation over the product of squared deviations. It now reads as the product of two separate
sums:

\begin{newtext}
\[
r = \frac{\sum(o_i-\mu_o)(f_i-\mu_f)}
         {\sqrt{\sum(o_i-\mu_o)^2\ \sum(f_i-\mu_f)^2}}
\]
\end{newtext}

We have checked the remaining equations; the KGE and PBIAS expressions are correct as
published. The reported correlation values were computed with a standard library function and
are unaffected by the typesetting error.

\subsection*{2.14 Nominal capacity}

\begin{rc}
(L107) 589 megawatts of ``nominal'' hydropower capacity. Is this installed capacity or
average generation? Please clarify.
\end{rc}

It is installed nameplate capacity. ``Nominal'' was ambiguous and we have replaced it. The
text now reads ``589 megawatts of installed nameplate hydropower capacity across 13
generating stations on the Connecticut and Deerfield Rivers in ISO New England,'' which also
makes clear that the system total covers many more facilities than the three evaluated here.

\subsection*{2.15 Relationship among the three sites}

\begin{rc}
(L109) Add a sentence describing how these 3 sites are related, given they are
downstream from one another one could assume they are operated jointly. Also good to note
why you chose these sites over all the other sites (green dots in the figure).
\end{rc}

Agreed, and this omission made the results harder to interpret than necessary. Your inference
is correct. The three dams lie in series on the Connecticut River mainstem, Wilder upstream,
then Bellows Falls, then Vernon, and are operated jointly, with releases from each becoming
inflow to the next. This is why the drainage areas are nested and why bias propagates
downstream, an effect we noted in the results without giving the reader the information needed
to understand it. We now quantify the propagation (comment 1.5).

On site selection, these three are the largest generators in the system, they are the dams
subject to the FERC forebay constraint that motivated the tailored metrics, and they are where
the operators concentrate their scheduling attention. The three tributaries were included
because they are the only significant local inflows, entering between Wilder and Bellows
Falls. The remaining facilities in Figure 2 are smaller and are not covered by the same
license conditions. Both points are now in the text.

\subsection*{2.16 Anonymization and reproducibility}

\begin{rc}
(L117) Anonymising the commercial forecast products is understandable but limits
reproducibility. Consider at minimum describing the general class of model used by each
(e.g., physics-based, statistical, machine learning) so readers can contextualise the
results.
\end{rc}

Agreed. Reviewers 1 and 4 asked for the same thing. Forecast A is a conceptual rainfall-runoff
model with hydraulic routing driven by commercial numerical weather prediction, issued daily;
forecast B is a machine learning model trained on gauge and remotely sensed observations,
issued twice daily; the in-house forecast combines RFC forecasts with routing and operator
judgment. Please see comment 1.4 for the full text and for what we cannot disclose.

We have also added the inability to inspect commercial products to the limitations, since it
constrains any evaluation of this kind and is part of what motivates judging products on
operational performance.

\subsection*{2.17 SI units}

\begin{rc}
(Figure 1) Y-axis units are ``cfs'' (cubic feet per second). Consider also providing SI
units for international readers.
\end{rc}

Added. The Figure 1 caption now gives the conversion (1,000 cfs $\approx$ 28.3 m$^3$/s), and
we give SI equivalents at first use in the text: the 0.5 foot forebay limit is noted as
approximately 15 cm, as you observed, and the largest observed inflow is given in both units.
We have kept cfs as the primary unit because it is the unit the operators and the FERC license
use, and the tailored metrics are expressed in feet of forebay for the same reason. We have
not added a secondary axis to the figure itself; if the editor prefers that we will do so.

\subsection*{2.18 Seasonal performance}

\begin{rc}
(L134) ``Forecasts tend to perform worse in the winter which is the wet season and worse
in the summer which is the dry season.'' This says forecasts perform worse in both wet and
dry seasons. Please clarify when forecasts actually perform well and why. The following
sentence about snow melt season partially addresses this but the logic is unclear.
\end{rc}

You are right that the sentence is self-contradicting and the logic was garbled. In
re-examining the monthly results we also found the original characterization of the seasons was
wrong: winter is not the wet season in this basin. Mean inflow at Vernon is highest in March
and April and lowest in September through November, so the high-flow season is spring, driven
by snowmelt, and that is also where skill is highest. We have rewritten the passage:

\begin{newtext}
``Skill is highest during the spring snowmelt period, roughly March through May, for all three
products. It is lowest in late summer and fall for the commercial products, with A showing the
sharpest degradation of any product in any season in January and February, where pooled
day-ahead KGE falls to 0.39 and $-0.40$ respectively. The seasons where performance is poorest
are those in which inflow is driven by precipitation events: rain and rain-on-snow in winter
and convective storms in summer. Both depend on precipitation forecasts which can be highly
uncertain. Snowmelt-driven inflow is comparatively predictable because it is governed by an
accumulated snowpack that can be observed in advance and melts over a period of weeks, so
skill depends more on temperature, which is forecast more reliably than precipitation.''
\end{newtext}

The original text also gave the snowmelt season as April through July, which does not match the
monthly inflow pattern; it is now March through May.

\subsection*{2.19 The inflow-equals-outflow assumption}

\begin{rc}
(L146) The 0.5 foot ($\approx$15 cm) operating range is an important constraint that
motivates a strong tailored metric. However, the forebay exceedance analysis assumes an
inflow=outflow condition which is operationally unrealistic as the authors note. It would
strengthen the analysis to discuss how operators actually deviate from this assumption and
what the real-world exceedance rates might be.
\end{rc}

A fair point. Operators do not hold pass-inflow conditions: they schedule generation to a
market bid, which means departing from inflow during the day and recovering the forebay
overnight, and they intervene when the forebay approaches the band. We have expanded the text
to say what the metric does and does not represent:

\begin{newtext}
``We note that the pass inflow assumption is not how the system is actually operated. In
practice operators schedule generation to a market bid, which typically means departing from
inflow during the day and recovering the forebay overnight. The metric is idealized and
intended to be an upper limit on the compliance risk a forecast carries rather than a
prediction of observed exceedance rates. Capturing the actual risk of exceedance given a
certain operation requires a scheduling model with direct operator oversight.''
\end{newtext}

We consider this the right quantity for comparing forecast products because it isolates
forecast quality from the operators' skill at correcting for it. Realized exceedance is less
frequent than these probabilities imply, but realized rates reflect both the forecast and the
effort spent compensating for it, and that effort is itself a cost of a poorer forecast.

We have not estimated real-world exceedance rates. Doing so credibly requires representing the
intervention behavior, which returns to the scheduling model coupling discussed under comment
4.9.

\subsection*{2.20 Extremes and the record length}

\begin{rc}
(L178--182) The discussion on limitations of real-time vs hindcast evaluation is good.
Consider also noting the limitation that almost two years of data may not capture the full
range of extreme events which are often the conditions where forecast quality matters most
for operational decisions.
\end{rc}

Agreed. Reviewer 4 asked in comment 4.2 what the period did contain, so we have addressed both
together, and checking the second question changed our answer to the first.

The period captured the July 2023 Vermont flood, which produced the largest calculated inflows
of the evaluation at Vernon and Bellows Falls, and a December 2023 rain-on-snow event that
produced the largest inflow at Wilder. Checking these against the long-term annual peak series
at the USGS reference gages places them in context: the July 2023 event ranks among the ten
largest flows on record on the Connecticut River main stem, but below the 1936, 1938, and 1927
floods and below the 2011 Tropical Storm Irene peak. The period also did not include a severe
multi-year drought. So the record covers a useful range but not the upper tail, and we now
report both the events observed and where they sit in the historical distribution. Please see
comment 4.2.

One point specific to this system is worth adding: the need to represent multiple high flow
events matters less at GRH than it might elsewhere, because operational flexibility is limited
during high flows and operators discount forecast errors in those conditions. The gap that
does matter for GRH is the absence of a strong multi-year drought, which we state as a
limitation.

\vspace{1em}
\hrule
\vspace{1em}

\section*{Reviewer 3}

\begin{rc}
However, proposed framework itself is not well described, such as what is new, different
from existing ones, and comparison of past studies are insufficient.
\end{rc}

Thank you for this review. Your observation that we implied the framework was novel without
describing the novelty was the most useful single comment we received, and the handling editor
highlighted it as well. Our response to the editor describes the new statement of contribution.
Your comments 3.4 and 3.5 have produced what we now consider the central result of the paper.

\subsection*{3.1 ``market prices''}

\begin{rc}
(l47) ``market prices'' --- Could you please clarify what it means?
\end{rc}

The sentence was too compressed. The point concerns why a hindcast evaluation is difficult for a
hydropower operator: reconstructing what a forecaster would have done in a past period requires
knowing the conditions they faced, and those conditions include the wholesale electricity prices
they were bidding against, which drive the generation schedule and hence release decisions.
Those price signals, and the operator's response to them, are not generally recoverable after
the fact. The revised passage is given under comment 2.6, which prompted a further correction to
the same sentences: the claim that a hindcast requires re-operating the system was overstated
for a quality assessment, and the requirement properly applies to a value assessment.

\subsection*{3.2 How feedback is reflected in the process}

\begin{rc}
(l50, \S2.1 real time inflow forecast evaluation) Step 1 through 3 are so general. I assume
that steps 4 and 5 are unique points in this manuscript. Please provide a more detailed
explanation of how feedback is reflected in step 4.
\end{rc}

You have identified the structure of the contribution correctly, and we now say so directly
rather than presenting all steps as equally novel. Steps 1 through 3 are generic and are
labeled as such.

Your comment also exposed an ambiguity: what was a single step for soliciting feedback
conflated gathering operator comment with actually changing the metric set. We have split them
into ``Gather feedback'' and ``Update metrics.'' The process is now six steps, with 4 through 6
being where it departs from retrospective verification.

We have expanded both with the mechanics of how feedback was collected and how it changed the
evaluation:

\begin{newtext}
``Feedback was collected in recurring meetings with the operators, held monthly for most of the
evaluation. Each meeting reviewed the current calculated metrics against recent operating
experience, and operators were asked to identify which recent forecast errors mattered
operationally, and which did not. The answers frequently diverged from what the generic metrics
indicated, especially during the high-flow season. Operators discounted large errors during
high-flow spill conditions, when they had little discretion over releases and forecast quality
was operationally irrelevant, while treating moderate errors with the potential to impact the
forebay compliance band as serious. These ongoing discussions motivated the tailored metrics:
the forebay exceedance metrics were designed in direct response to operators stating that
compliance risk, rather than average error, determined whether they trusted a forecast.''
\end{newtext}

We have also added a table tracing each tailored metric to the operator concern that prompted it
and the approximate point at which it was added, so the process is auditable rather than
asserted:

\begin{center}
\begin{tabular}{p{0.24\linewidth}p{0.42\linewidth}l}
\toprule
Metric & Operator concern & Added \\
\midrule
Day-ahead volumetric skill & Day-ahead market bids are set from next-day inflow volume, not
hourly shape & Months 1--2 \\
Seasonal metric breakdown & Products were suspected to behave differently in snowmelt and
rain driven regimes & Months 3--4 \\
Cumulative forebay error & Average error does not indicate whether the forebay band would be
breached & Months 4--6 \\
Forebay exceedance from a non-neutral start & Operators rarely begin a day at mid-band, which
is the harder and more common case & Months 5--6 \\
\bottomrule
\end{tabular}
\end{center}

We should be candid that the months are approximate. The evaluation predates this paper and we
did not keep a formal log of metric provenance, so these are reconstructed from meeting
materials and correspondence. The sequence and the motivations are accurate, and we have
described the timing as approximate rather than implying precision we cannot support. A
practical recommendation follows, which we now state: an evaluation of this kind should log
metric provenance as it goes, since the record of why each metric exists is part of what makes
the results defensible to the operators later.

\subsection*{3.3 Is a one-year learning period required?}

\begin{rc}
(l62, ``5. Iterate'') This implies that at least one year learning period is required. Is
this understanding correct?
\end{rc}

Your understanding is correct for a full evaluation, and your question identified that the
original wording conflated two durations. We have separated them:

\begin{newtext}
``The iteration in step 6 involves two timescales that should not be conflated. Metric
development was completed within roughly the first six months of our evaluation, after which
the iteration served to accumulate data rather than to identify new metrics. A minimum
defensible skill assessment requires at least one full annual cycle, since forecast performance
may vary substantially by season and different forecast models may rank differently in each
season. Operators should therefore expect the set of metrics to stabilize well before the
evaluation period ends.''
\end{newtext}

Our own results bear on the second point. Pooled day-ahead KGE for forecast A ranges from
$-0.40$ in February to 0.88 in May, so an evaluation conducted in spring alone would have
reached a substantially more favorable conclusion about that product than the full record
supports. We have also noted the tension the operators raised, that industry preference runs
toward shorter evaluations, sometimes a single season, while a single season samples only one
hydrologic outcome. The framework does not resolve this, and we recommend that a real-time
evaluation be paired with a retrospective check where one is feasible.

\subsection*{3.4 Comparison with traditional methods}

\begin{rc}
(l105, result of case study) Were all the results derived through the proposed framework? If
so, please compare those result with those obtained via ``traditional'' methods. This
conmparison may enable to enhance the effectiveness of the proposed framework.
\end{rc}

All results were produced through the framework, and we agree the comparison you describe is
what demonstrates its effectiveness. We have added it in two forms.

First, a discussion subsection, ``Comparison with retrospective forecast evaluation.'' A
retrospective evaluation provides larger sample sizes, a wider range of hydrologic conditions,
and better coverage of extremes, since hindcasts can span decades where our evaluation spans
less than two years. Against that, translating thousands of historical forecasts into
operational decisions may be prohibitively expensive, and even where historical operations could
be reconstructed it is difficult or impossible to reconstruct the exact market conditions at the
time of a decision. A real-time evaluation reflects actual operational conditions including
outages and market feedbacks, and requires no additional data or extra operator effort. We are
explicit that the real-time framework is not a replacement for retrospective verification where
a hindcast is feasible; it is a practical alternative where one is not.

Second, the comparison of conclusions. A conventional verification of these forecasts, computing
NSE, KGE, PBIAS, and correlation over the full record, concludes that forecast B outperforms
forecast A and is comparable to the in-house forecast. That conclusion is correct and this
framework reproduces it. What the conventional approach omits is the compliance risk finding
that carried the most weight with the operators, detailed under comment 3.5. We are careful not
to overstate the point: the two approaches agree on the ranking, and what the tailored metric
adds is a dimension of comparison the generic metrics leave invisible rather than a
contradiction of them.

\subsection*{3.5 Advantage of tailored metrics}

\begin{rc}
(l139, \S3.1 Tailored metrics) Please explain the advantage of introducing tailored metrics.
For example, it would be helpful to show a case where using tailored metrics led to a
different forecast onclusion compared to general metrics.
\end{rc}

This was the right thing to ask for and we have added that case as a result rather than leaving
it for the reader to infer.

We should be precise about the strength of the claim. The generic and tailored metrics do not
reach opposite conclusions about which product is better; forecast B outperforms forecast A on
both, so the overall ranking is consistent. What the tailored metric does is separate the
products in a dimension the generic metrics leave invisible, and quantify the separation in
units the operators are accountable for.

The forebay exceedance metric is that case. From a neutral forebay starting position, the
probability of breaching the compliance band is below 10\% for every product at every lead time
in the 1--3 day range, and on that basis the products are difficult to distinguish. From the
non-neutral starting position that operators are routinely in, they separate sharply: at Bellows
Falls forecast A reaches 23.7\% on day 1 and 11.9\% on day 2, against 1.8\% and 2.4\% for
forecast B; at Vernon forecast A reaches 10.1\% by day 3 against 5.2\%. A roughly tenfold
difference in compliance risk at Bellows Falls is not visible in the NSE, KGE, PBIAS, or
correlation values, none of which know that a penalty attaches to a particular forebay
threshold.

The non-neutral case exists in the analysis only because the operators asked for it. The
neutral-start case, which is what an analyst would naturally compute and what our original draft
emphasized, shows nothing. This is the mechanism by which operator participation in metric
design changes the result rather than merely improving its reception.

We have added the general point to the manuscript:

\begin{newtext}
``The general point we would like to convey in this section is the importance of computing an
array of traditional and tailored metrics. Hydrologic forecasts are highly complex and nuanced,
and rarely will there be a cut and dry answer to which is better. Tailored metrics provide
information to operators in ways that might be obscured by traditional evaluation metrics,
ultimately providing more robust information for investment decisions.''
\end{newtext}

On comparison to past studies, the introduction already discussed the head-to-head evaluation
exercises (the HEPEX community experiments and the Bureau of Reclamation forecast rodeos) and we
have sharpened the statement of what distinguishes this work: those exercises evaluate forecast
quality in general, while this framework evaluates fitness for one utility's operating
constraints under live conditions.

\vspace{1em}
\hrule
\vspace{1em}

\section*{Reviewer 4}

\begin{rc}
The overall idea is relevant and practically useful for hydropower operators seeking to assess
competing forecast products. The paper is generally well-written, and the iterative,
operator-inclusive evaluation approach adds value.
\end{rc}

Thank you for the careful reading and for the specific, actionable comments. We have addressed
all nine. Your comment 4.2 led us to check the evaluation period against the long-term record,
which corrected a claim we would otherwise have overstated.

\subsection*{4.1 Smoothing window size and sensitivity}

\begin{rc}
Lines 80--81 recommend a triangular smoother for calculated inflows but do not specify the
smoothing window size (number of days or timesteps). This is a critical parameter because the
choice of window directly affects computed evaluation metrics. The window size must be stated
explicitly, and a sensitivity analysis or justification for the chosen value would strengthen
the paper.
\end{rc}

You are right that this parameter was left unstated and that it affects the metrics. We have
added both the value and a sensitivity analysis. Reviewer 2 raised the same point in comment 2.8.

The smoothing is a 24-hour triangular window on hourly calculated inflow, implemented as two
successive 24-hour moving averages, which produces the triangular weighting. The window was
matched to the operational timescale: the day-ahead volumetric forecast is the quantity of
interest, so smoothing at 24 hours removes sub-daily forebay noise without attenuating the daily
signal being evaluated.

We repeated the day-ahead evaluation at window widths of 0 (no smoothing), 6, 12, and 24 hours.
Day-ahead NSE at each dam:

\begin{center}
\begin{tabular}{llrrrr}
\toprule
Location & Forecast & 0 h & 6 h & 12 h & 24 h \\
\midrule
Wilder        & A   & 0.70 & 0.72 & 0.72 & 0.72 \\
              & B   & 0.82 & 0.83 & 0.83 & 0.83 \\
              & GRH & 0.87 & 0.87 & 0.88 & 0.88 \\
Bellows Falls & A   & 0.50 & 0.51 & 0.52 & 0.55 \\
              & B   & 0.90 & 0.90 & 0.90 & 0.91 \\
              & GRH & 0.84 & 0.83 & 0.84 & 0.85 \\
Vernon        & A   & 0.79 & 0.79 & 0.79 & 0.81 \\
              & B   & 0.88 & 0.88 & 0.88 & 0.89 \\
              & GRH & 0.88 & 0.88 & 0.88 & 0.89 \\
\bottomrule
\end{tabular}
\end{center}

Metrics improve modestly with window width, as expected since smoothing removes observation
noise that no forecast can reproduce. The largest change is 0.05 NSE (Bellows Falls, forecast A).
The improvement is not strictly monotonic: the in-house forecast at Bellows Falls dips from 0.84
to 0.83 between 0 and 6 hours before recovering, which we attribute to sampling rather than to
any property of the smoother. The ranking of the three products is unchanged at every window
width at every location, and the gap between the two commercial products is stable: forecast B
leads forecast A by 0.39 NSE at Bellows Falls with no smoothing and by 0.37 with 24-hour
smoothing. The 24-hour column reproduces Table 1 exactly, so the sensitivity analysis and the
main results come from the same pipeline.

We have added the table and the following caution, which is the more generally useful part of the
answer:

\begin{newtext}
``Evaluators should nonetheless state the window explicitly and check that their conclusions are
insensitive to it, since a window wide enough to attenuate the signal of interest would flatter
forecasts that under-predict variability. We recommend matching the window to the decision
timescale rather than choosing the window that maximizes metric values.''
\end{newtext}

\subsection*{4.2 Hydrologic representativeness of the evaluation period}

\begin{rc}
Line 116 states the evaluation ran from July 2023 to May 2025. The authors should discuss whether
this period captured any notable extreme hydrologic events (major floods, drought conditions) and
how climatologically representative it is relative to the long-term record for the Connecticut
River basin. This context is important for understanding the generalizability of the approach.
\end{rc}

This was a valuable comment and answering it properly changed two things in the manuscript.

First, the stated period was wrong. The record analyzed runs from June 2023 to April 2025, and the
three products do not cover it uniformly: forecast B spans the full period (657 day-ahead
forecasts per dam) while forecast A and the in-house forecast both end in November 2024 (539 and
528 respectively). We have corrected the period, added the per-product sample sizes to Table 1,
and verified that restricting the comparison to the 510 days on which all three products are
present leaves the ranking unchanged at all three dams, with individual metrics shifting by no
more than 0.03 NSE. We now report this partial coverage as a characteristic feature of real-time
evaluation rather than smoothing over it.

Second, on the events. The period captured the July 2023 Vermont flood, which produced the largest
calculated inflows of the evaluation at Vernon (83,100 cfs, 2,353 m$^3$/s) and Bellows Falls
(67,000 cfs), and a December 2023 rain-on-snow event that produced the largest inflow at Wilder
(27,600 cfs) and the second largest at Vernon and Bellows Falls. It also included spring freshets
in 2024 and 2025 and dry summer conditions in both 2023 and 2024. We checked the July 2023 event
against the long-term annual peak series at the USGS reference gages: it ranks among the ten
largest flows on record on the Connecticut River main stem, though below the 1936, 1938, and 1927
floods and below the 2011 Tropical Storm Irene peak. So the period was not hydrologically quiet,
but neither did it sample the upper tail. We have added a subsection reporting this and the
following:

\begin{newtext}
``The period therefore sampled a reasonable range of the conditions relevant to GRH operations.
It did not, however, sample the full distribution of hydrologic conditions. Notably, the July 2023
event produced one of the ten largest flows ever recorded on the main stem of the Connecticut
River. The evaluation also did not include any severe multi-year drought. Any real-time evaluation
is bounded in this way, and the conclusions below should be taken as conditional on the observed
period.''
\end{newtext}

Reviewer 2 asked in comment 2.20 for the limitation to be stated explicitly, which we have added
to the limitations discussion. One caveat on the comparison is worth noting: the gage annual peak
series are instantaneous peak discharge at the gage, while our values are 24-hour smoothed
calculated inflow at the dams, so the two are not strictly commensurable. We have therefore phrased
the manuscript statement as an approximate ranking rather than a numerical comparison.

\subsection*{4.3 Characterization of the commercial forecasts}

\begin{rc}
Line 117 The authors anonymize the two commercial forecasts as A and B. While commercial
confidentiality is understandable, the paper should provide at least a general characterization of
each forecast's methodology (e.g., physics-based, machine-learning-based, hybrid, statistical).
Without this, readers cannot interpret the results in a broader methodological context or assess
the generalizability of the findings.
\end{rc}

Agreed, and Reviewers 1 and 2 raised the same point. Please see comment 1.4 for the characterization
we have added and for what we cannot disclose. In summary: forecast A is a conceptual
rainfall-runoff model with hydraulic routing driven by commercial numerical weather prediction,
issued once daily out to roughly 8 days; forecast B is a machine learning model trained on gauge
and remotely sensed observations, issued twice daily out to 10 days; the in-house forecast combines
RFC forecasts with routing and operator judgment and extends about 32 hours. The issue frequencies
and horizons are taken from the forecast archive rather than vendor documentation.

We have also added that neither vendor released model internals or calibration diagnostics, which
is relevant to the framework's motivation rather than merely a caveat.

\subsection*{4.4 Extension to probabilistic and ensemble forecasts}

\begin{rc}
The paper focuses only on deterministic forecast evaluation (lines 42--49). While this is
acknowledged, the framework could be briefly extended to probabilistic or ensemble forecasts using
metrics such as CRPS, reliability diagrams, or rank histograms. This would make the description
more complete and forward-looking, especially given the authors' own statement that probabilistic
forecasts are gaining in popularity.
\end{rc}

We agree and have added a methods subsection. Reviewer 2 asked for the same in comment 2.5:

\begin{newtext}
``The framework described above is agnostic to whether the forecast being evaluated is
deterministic or probabilistic. Probabilistic forecasts are valuable for assessing the risk
associated with certain operational decisions, which is particularly useful for compliance or
market related concerns. Our evaluation framework can be used with either probabilistic or
deterministic forecasts. The main difference being in the choice of metrics and visualizations, the
overall process remains the same. Where ensemble or probabilistic products are available, metrics
such as the continuous ranked probability score (CRPS) can be used. CRPS has the useful property
that it reduces to mean absolute error (MAE) for a deterministic forecast, so the two can be
compared on a common scale. Diagrams such as rank histograms and reliability diagrams are useful
for assessing properties of ensemble forecasts.''
\end{newtext}

Two connections are worth drawing out. First, the forebay exceedance metric we develop is already a
probabilistic statement derived from a deterministic forecast, computed across the empirical
distribution of that forecast's past errors; given an ensemble it would be computed directly from
the members, which is both simpler and better founded. Second, operators bidding into a day-ahead
market must submit a single quantity, so a probabilistic forecast must be reduced to a point value
at the moment of the bid, which is why improvements in deterministic skill remain directly useful.
We retain that point in the limitations.

This is a description of how the framework would extend rather than a demonstration, since all
three products evaluated here are deterministic and we have no ensemble product to apply it to.

\subsection*{4.5 Hyphenation of ``real time''}

\begin{rc}
(Introduction) Throughout the manuscript: ``real time'' should be hyphenated when used as a
compound adjective (e.g., ``real-time evaluation''). Usage is inconsistent.
\end{rc}

Corrected throughout, including the title. Reviewer 2 raised the same point in comment 2.3, where we
describe the treatment. In brief: the manuscript contained 29 instances of ``real time'' and none
hyphenated; we reviewed each individually, hyphenating in attributive position and leaving the
adverbial phrase (``used in real time'') unhyphenated.

\subsection*{4.6 KGE citation}

\begin{rc}
(Methods, line 92) Remove the placeholder ``(refs)'' and insert the actual KGE citation (Gupta et
al.\ 2009).
\end{rc}

Corrected. The placeholder is replaced with Gupta et al.\ (2009), and we have added Nash and
Sutcliffe (1970) for NSE in the same sentence. We also corrected the misspelling Reviewer 2 caught
in comment 2.9: the manuscript read ``King-Gupta Efficiency'' and now reads ``Kling-Gupta
Efficiency.''

While checking this we found and fixed a broken citation key elsewhere: the introduction cited
\texttt{kratzertLearningUniversalRegional2019} where the bibliography entry is
\texttt{kratzertLearningUniversalRegional2019a}, so that citation was not resolving. All citation
keys now resolve against the bibliography.

\subsection*{4.7 RMSE or MAE in Table 1}

\begin{rc}
(Results and Case Study, Table 1) Consider adding RMSE or MAE values alongside the existing metrics
to provide a sense of forecast error magnitude in physical units.
\end{rc}

Agreed, and this was a real omission since every metric in the original table was dimensionless. We
have added RMSE in cfs to Table 1. The values are informative in themselves: forecast A's day-ahead
RMSE at Bellows Falls is 5,904 cfs against 2,615 for forecast B, which conveys the practical size of
the difference more directly than the NSE values of 0.55 and 0.91. Table 1 also gains the KGE skill
score column described under comment 1.6 and the per-product sample size described under comment 4.2.

\subsection*{4.8 Figure 5 caption}

\begin{rc}
(Figure 5 caption) The caption does not specify that the error bars represent the 2.5th and 97.5th
percentiles (as stated in the main text at lines 148--149). Figure caption needs to be
self-contained.
\end{rc}

Corrected. The Figure 5 caption now states that points show the median across all forecasts and error
bars span the 2.5th to 97.5th percentile, and specifies the pass-inflow assumption and the
$\pm$0.5 foot compliance band. We reviewed all captions for self-containment as part of the caption
revision described under comment 1.8.

\subsection*{4.9 What coupling to scheduling models would entail}

\begin{rc}
(Discussion, lines 173--174) The authors note that coupling this evaluation with hydropower
scheduling models is an important next step. This point could be strengthened by briefly discussing
what such a coupling would entail and what data would be needed.
\end{rc}

Agreed, the original text asserted the next step without describing it. We have expanded the passage:

\begin{newtext}
``One difficulty in this approach is that scheduling models may be configured to reproduce a
utility's current bidding strategy, which was itself developed around the accuracy of its existing
forecast. A more accurate forecast may show little revenue benefit simply because the model is not
set up to exploit it. Quantifying the value of improved forecasts may therefore require revising the
scheduling model and the operating strategy alongside any new forecast.''
\end{newtext}

On the data required, such a coupling would route each candidate forecast through a scheduling model
to produce the generation schedule it implies, then value that schedule against realized inflows and
market prices, so that forecast differences are expressed in revenue rather than skill. This needs
the utility's scheduling model with its operating constraints and unit efficiency curves, historical
day-ahead and real-time market prices at the relevant node, the realized inflow record, and the
archive of forecasts as issued so that schedules are constructed only from information available at
bid time. That last requirement is a further reason a real-time evaluation is a useful precursor,
since it produces exactly this archive as a byproduct.

This connects to Reviewer 2's first major comment, since the same coupling is what a full forecast
value assessment would require, and to comment 2.19, since estimating realized forebay exceedance
rates needs the same machinery.

\vspace{1em}
\hrule
\vspace{1em}

\section*{Summary of changes}

\subsection*{Framing and positioning}

\begin{itemize}[leftmargin=1.4em,itemsep=0.15em]
\item New statement of contribution at the end of the introduction, distinguishing the framework from
retrospective verification on two grounds: operator participation in metric selection and design, and
evaluation on the operational data feed (editor, 1.1, 2.2, 3 general).
\item Introduction gap statement rewritten so the gap is the absence of an evaluation procedure,
stating why hindcast verification is often infeasible and why archived operational forecasts are a
poor substitute for continuously updated, operator-adjusted, or third-party products (2.4).
\item Process restructured from five steps to six, separating ``Gather feedback'' from ``Update
metrics.'' Steps 1--3 labeled as generic practice, with novelty claimed only for 4--6 (2.2, 3.2).
\item New discussion subsection ``Comparison with retrospective forecast evaluation'' (3.4).
\item New discussion subsection ``From forecast quality to forecast value,'' citing Farkas et al.\
(2025) and giving operators a two-stage pathway to a value assessment (2.1).
\item Abstract revised to state what distinguishes the framework and to reference the case where a
tailored metric separated products that the standard metrics could not.
\end{itemize}

\subsection*{New results and analysis}

\begin{itemize}[leftmargin=1.4em,itemsep=0.15em]
\item New result: the forebay exceedance metric separates the two commercial products where the
generic metrics do not. From a neutral start all products are below 10\% exceedance; from the
non-neutral start operators are routinely in, forecast A reaches 23.7\% at Bellows Falls against
1.8\% for forecast B (3.5).
\item KGE skill score against the persistence forecast added as a column to Table 1, with a new
equation defining a skill score generally. The skill score reverses the apparent cross-site ordering
of forecast A, which is weakest at Vernon (0.47) rather than at Bellows Falls (0.65) as the raw NSE
and KGE values suggest (1.6, 4.7).
\item Smoothing window stated explicitly (24-hour triangular) with a new sensitivity table over 0, 6,
12, and 24 hours showing the ranking is unchanged; the 24-hour column reproduces Table 1 exactly
(4.1, 2.8).
\item RMSE in cfs added to Table 1 (4.7).
\item Per-product sample sizes added to Table 1; common-period check reported, showing the ranking is
unchanged (4.2).
\item Downstream bias propagation quantified rather than asserted (1.5, 2.15).
\item Evaluation period corrected to June 2023--April 2025 and per-product coverage reported (4.2).
\item New subsection on the hydrology of the evaluation period, including the July 2023 Vermont flood
and December 2023 rain-on-snow event, with their position in the long-term annual peak record (4.2,
2.20).
\item Seasonal passage corrected: spring is the high-flow season and where skill is highest, not
winter; the self-contradiction removed and the mechanism explained (2.18).
\end{itemize}

\subsection*{New descriptive material}

\begin{itemize}[leftmargin=1.4em,itemsep=0.15em]
\item New subsection on GRH operating context: run-of-river operation, the daily cycle, the day-ahead
market, and the FERC forebay limit (1.3).
\item New subsection on the forecast products, characterizing all three by model type, forcing, issue
frequency, and horizon, with the inability to inspect commercial products noted (1.4, 2.16, 4.3).
\item New table of drainage areas and USGS reference gages for all six evaluation points, plus basin
description including mainstem nesting and upstream flood control regulation (1.5).
\item New table tracing each tailored metric to the operator concern that prompted it and the
approximate month it was added (3.2).
\item Steps 4 and 5 expanded with the mechanics of feedback collection and how it changed the metric
set (3.2).
\item Iteration step clarified to separate metric convergence from record length (3.3).
\item New methods subsection extending the framework to probabilistic and ensemble forecasts (CRPS,
rank histograms, reliability diagrams) (4.4, 2.5).
\item Expanded discussion of what coupling to a scheduling model would entail and why the current
bidding strategy limits measurable benefit (4.9).
\item Expanded treatment of the pass-inflow assumption: the metric measures unmitigated compliance
risk, realized rates are lower, and estimating them requires the scheduling model coupling (2.19).
\item Added that calculated inflow is itself operations-dependent, so the reference data is not
independent of operations (2.6).
\item Added that operational reliability is part of what a real-time evaluation measures, since a
forecast that cannot be produced operationally should not be considered (2.2).
\end{itemize}

\subsection*{Corrections}

\begin{itemize}[leftmargin=1.4em,itemsep=0.15em]
\item Pearson correlation equation corrected: the denominator now has separate summations over each
squared deviation term (2.13).
\item Bias definition corrected to distinguish error/residual from systematic bias, and made
consistent in sign with the PBIAS equation (2.12).
\item KGE rationale corrected: the advantage is decomposition into correlation, variability, and bias
ratios, not insensitivity to large values (2.10).
\item ``King-Gupta'' corrected to ``Kling-Gupta''; the ``(refs)'' placeholder replaced with Gupta et
al.\ (2009) and Nash and Sutcliffe (1970) (2.9, 4.6).
\item Knoben et al.\ (2019) citation moved adjacent to the $-0.41$ threshold claim it supports, and
the consequence stated (2.11).
\item Broken citation key \texttt{kratzertLearningUniversalRegional2019} fixed; all keys verified
against the bibliography (4.6).
\item Persistence redescribed as ``a simple but competitive'' forecast rather than the worst possible
forecast, and its meaningful horizon stated (2.7, 1.2).
\item ``Nominal'' capacity clarified as installed nameplate capacity, with the number of generating
stations given (2.14).
\item Added that the three dams are in series on the mainstem and operated jointly, and why these six
sites were selected (2.15).
\item ``real time'' hyphenated in attributive position throughout, including the title (2.3, 4.5).
\item SI units given at first use of cfs and for the 0.5 foot forebay limit, and in the Figure 1
caption (2.17).
\item All figure captions made self-contained, including the Figure 5 percentile definition, the
no-skill threshold values, and the qualification on persistence at long lead times (1.8, 4.8, 1.2).
\item Tables rebuilt with grouped rows, consistent significant figures, and units in headers (1.7).
\item Farkas et al.\ (2025) added to the bibliography for the forecast value discussion (2.1).
\end{itemize}

\subsection*{Requests we were not able to fully satisfy}

We would rather name these than let them pass as addressed.

\begin{itemize}[leftmargin=1.4em,itemsep=0.15em]
\item \textbf{Climatological benchmark in Figures 3 and 4} (1.2). We recommend one in the text but have
not added one, because a defensible daily climatology of calculated inflow at the dams requires a long
record that begins only with the evaluation, and a gage-derived substitute would not be commensurable
with the quantity verified.
\item \textbf{Calibration and validation diagnostics for the commercial products} (1.4). Not released
to us by either vendor. We state the gap and treat it as a finding about the procurement situation.
\item \textbf{Redrawing the figures} (1.8). Captions are substantially revised and we have described
what we would change, but we have not regenerated the figures and would welcome specific guidance on
which aspects are considered deficient.
\item \textbf{A full forecast value assessment} (2.1). We cite the literature and give a two-stage
pathway but compute no value metric, which would require the utility's damage function and a model of
operator response. We have chosen a review citation over a detailed comparison of specific value
metrics, and are glad to expand that treatment if preferred.
\item \textbf{Realized forebay exceedance rates} (2.19). Requires representing intervention behavior,
which needs the scheduling model coupling described under 4.9.
\item \textbf{Demonstrated probabilistic evaluation} (2.5, 4.4). Described but not demonstrated, as all
three products evaluated are deterministic.
\end{itemize}

\vspace{1em}

We thank the reviewers again for the care they took. In three places their questions caused us to find
and correct errors that would otherwise have gone to print: the stated evaluation period, the seasonal
characterization, and the cross-site comparison of NSE values. Reviewer 1's request that we check the
NSE values at Bellows Falls (comment 1.6) was the most productive comment we received, since following
it up led to the skill score column and to a clearer treatment of why raw efficiency metrics should not
be compared across locations.

\end{document}
